\documentclass[11pt, a4paper]{article}

% bfw-style.tex — gemeinsame Style-Definitionen für alle Handouts/Aufgabenhefte
% Voraussetzung: Kompilation mit xelatex (wegen fontspec)

\usepackage[ngerman]{babel}
\usepackage{geometry}
\geometry{a4paper, top=2.2cm, bottom=2.2cm, left=2.3cm, right=2.3cm}

\usepackage{fontspec}
% Fallback-Kette: Source Sans Pro -> Calibri -> Segoe UI -> Latin Modern Sans
% (Latin Modern ist immer mit MiKTeX vorhanden)
\IfFontExistsTF{Source Sans Pro}{%
  \setmainfont{Source Sans Pro}[Ligatures=TeX]%
}{\IfFontExistsTF{Calibri}{%
  \setmainfont{Calibri}[Ligatures=TeX]%
}{\IfFontExistsTF{Segoe UI}{%
  \setmainfont{Segoe UI}[Ligatures=TeX]%
}{%
  \setmainfont{Latin Modern Sans}[Ligatures=TeX]%
}}}
\IfFontExistsTF{Source Code Pro}{%
  \setmonofont{Source Code Pro}[Scale=0.92]%
}{\IfFontExistsTF{Consolas}{%
  \setmonofont{Consolas}[Scale=0.92]%
}{%
  \setmonofont{Latin Modern Mono}[Scale=0.92]%
}}

\usepackage{xcolor}
% Farbpalette: hell, freundlich, modern
\definecolor{bfwPrimary}{HTML}{0EA5A4}    % Petrol/Türkis - Primär
\definecolor{bfwPrimaryDark}{HTML}{0F766E} % dunkler Petrol für Headings
\definecolor{bfwAccent}{HTML}{F59E0B}     % Amber - Akzent
\definecolor{bfwSuccess}{HTML}{10B981}    % Grün - Erfolg / Lösung
\definecolor{bfwWarn}{HTML}{F97316}       % Orange - Achtung
\definecolor{bfwInk}{HTML}{0F172A}        % fast Schwarz
\definecolor{bfwInkSoft}{HTML}{475569}    % gedämpftes Grau
\definecolor{bfwBgSoft}{HTML}{F8FAFC}     % off-white
\definecolor{bfwBgTeal}{HTML}{ECFEFF}     % helles Türkis
\definecolor{bfwBgAmber}{HTML}{FEF3C7}    % helles Gelb
\definecolor{bfwBgRose}{HTML}{FFE4E6}     % helles Rosé für Eselsbrücken

\usepackage[colorlinks=true, linkcolor=bfwPrimaryDark, urlcolor=bfwPrimaryDark]{hyperref}
\usepackage{enumitem}
\setlist[itemize]{leftmargin=*, itemsep=2pt, topsep=2pt}
\setlist[enumerate]{leftmargin=*, itemsep=2pt, topsep=2pt}

\usepackage[most]{tcolorbox}
\usepackage{booktabs}
\usepackage{tikz}
\usetikzlibrary{decorations.pathmorphing, positioning, shapes.geometric, arrows.meta}

% Merkkästchen — wichtige Definition / Kernpunkt
\newtcolorbox{merksatz}[1][]{%
  enhanced, breakable,
  colback=bfwBgTeal,
  colframe=bfwPrimary,
  boxrule=0.6pt,
  arc=4pt,
  left=10pt, right=10pt, top=8pt, bottom=8pt,
  fonttitle=\bfseries\color{bfwPrimaryDark},
  title={\faLightbulb\ Merksatz},
  #1
}

% Eselsbrücke — Mnemoniks
\newtcolorbox{eselsbruecke}[1][]{%
  enhanced, breakable,
  colback=bfwBgRose,
  colframe=bfwAccent,
  boxrule=0.6pt,
  arc=4pt,
  left=10pt, right=10pt, top=8pt, bottom=8pt,
  fonttitle=\bfseries\color{bfwWarn},
  title={Eselsbrücke},
  #1
}

% Beispiel-Box
\newtcolorbox{beispiel}[1][]{%
  enhanced, breakable,
  colback=bfwBgSoft,
  colframe=bfwInkSoft,
  boxrule=0.4pt,
  arc=3pt,
  left=10pt, right=10pt, top=8pt, bottom=8pt,
  fonttitle=\bfseries\color{bfwInk},
  title={Beispiel},
  #1
}

% Lösungs-Box (für Aufgabenhefte)
\newtcolorbox{loesung}[1][]{%
  enhanced, breakable,
  colback=bfwBgAmber!40,
  colframe=bfwSuccess,
  boxrule=0.6pt,
  arc=3pt,
  left=10pt, right=10pt, top=8pt, bottom=8pt,
  fonttitle=\bfseries\color{bfwSuccess!70!black},
  title={Lösung},
  #1
}

% Glossar-Eintrag
\newcommand{\glossareintrag}[2]{%
  \par\noindent\textbf{\color{bfwPrimaryDark}#1} \enspace #2\par\smallskip
}

% Section-Stil — etwas farbiger als Standard
\usepackage{titlesec}
\titleformat{\section}
  {\Large\bfseries\color{bfwPrimaryDark}}
  {\thesection}{0.6em}{}
\titleformat{\subsection}
  {\large\bfseries\color{bfwPrimary}}
  {\thesubsection}{0.5em}{}
\titleformat{\subsubsection}
  {\normalsize\bfseries\color{bfwInk}}
  {\thesubsubsection}{0.4em}{}

% TikZ-Stil "skizzenhaft" — etwas weicher als Rough.js, aber stabil
\tikzset{
  roughbox/.style = {
    draw=bfwPrimaryDark, thick, fill=bfwBgTeal,
    rounded corners=4pt, inner sep=6pt
  },
  roughaccent/.style = {
    draw=bfwAccent, thick, fill=bfwBgAmber,
    rounded corners=4pt, inner sep=6pt
  },
  roughline/.style = {
    draw=bfwInkSoft, thick, -{Stealth[length=2.5mm]}
  }
}

% Bullet-Symbole (bewusst kein FontAwesome — vermeidet Font-Installations-Probleme)
\newcommand{\faLightbulb}{$\bullet$}
\newcommand{\faBookmark}{$\bullet$}

% Header/Footer
\usepackage{fancyhdr}
\pagestyle{fancy}
\fancyhf{}
\renewcommand{\headrulewidth}{0pt}
\fancyfoot[L]{\small\color{bfwInkSoft}\thepage}
\fancyfoot[R]{\small\color{bfwInkSoft} BFW M-064 Beschaffung im Büromanagement}


\title{\vspace*{-1.5cm}\Huge\bfseries\color{bfwPrimaryDark}Tag~3: Bestellungen und Kaufverträge\\[0.4em]
  \large\color{bfwInkSoft}\textnormal{Rechtliche Grundlagen, Vertragsschluss und AGB}}
\author{\textsf{Beschaffung im Büromanagement}\\
\textsf{\small\copyright\ Can Siebert\,\textperiodcentered\,2026}}
\date{13.05.2026}

\begin{document}
\maketitle
\thispagestyle{empty}

\vspace{1em}
\begin{merksatz}[title={\bfseries Worum geht es heute?}]
Nachdem Sie an den letzten Tagen den Beschaffungsprozess organisatorisch durchlaufen
haben, wenden wir uns heute dem \emph{rechtlichen Rückgrat} jedes Beschaffungs­vorgangs
zu: dem Kaufvertrag. Wer im Beruf bestellt, schließt jeden Tag Verträge --- meistens
ohne es bewusst zu merken. Dieses Handout vermittelt die rechtlichen Grundlagen, die
Sie brauchen, um Verträge sicher abzuschließen, AGB richtig einzuordnen und im Streitfall
zu wissen, welche Rechte und Pflichten gelten.
\end{merksatz}

\tableofcontents
\vspace{1em}
\hrule
\vspace{1em}

\section{Rechtsgebiete: privat und öffentlich}

Das deutsche Recht gliedert sich grob in zwei große Gebiete. Das \emph{private Recht}
regelt die Beziehungen zwischen gleichberechtigten Personen und Unternehmen --- also
zwischen Käufer und Verkäufer, Mieter und Vermieter, Arbeitnehmer und Arbeitgeber.
Wichtigste Gesetze sind das Bürgerliche Gesetzbuch (BGB) und das Handelsgesetzbuch (HGB).
Das \emph{öffentliche Recht} regelt das Verhältnis zwischen Bürger und Staat ---
Über- und Unterordnung. Hierzu gehören Grundgesetz, Verwaltungsrecht, Steuerrecht und
Strafrecht.

\begin{merksatz}
Beschaffungs­verträge fallen ausnahmslos ins \textbf{private Recht}. Das BGB ist Ihre
wichtigste Rechtsquelle im Berufsalltag --- besonders die §§\,433 ff.\ (Kaufvertrag) und
§§\,305 ff.\ (AGB).
\end{merksatz}

\section{Rechtssubjekte und Geschäftsfähigkeit}

\subsection{Wer ist „Person" im Recht?}

Nur \emph{Rechtssubjekte} können Träger von Rechten und Pflichten sein und damit
Verträge abschließen. Man unterscheidet:

\textbf{Natürliche Personen} sind alle Menschen --- vom Tag der Geburt (§\,1 BGB) bis
zum Tod. Jeder Mensch ist also rechtsfähig.

\textbf{Juristische Personen} sind künstliche Rechts­personen, die das Gesetz mit
eigener Rechtspersönlichkeit ausstattet. Beispiele: Gesellschaft mit beschränkter
Haftung (GmbH), Aktiengesellschaft (AG), eingetragener Verein (e.\,V.), Städte und
Gemeinden. Sie können selbst Verträge schließen, klagen und verklagt werden.

\subsection{Geschäftsfähigkeit --- drei Stufen}

Rechtsfähig zu sein, heißt aber noch nicht, eigenständig wirksame Verträge schließen
zu können. Dafür braucht es \emph{Geschäfts­fähigkeit}. Das BGB unterscheidet drei
Stufen.

\textbf{Geschäftsunfähig} sind Kinder unter 7~Jahren (§\,104 Nr.\,1 BGB) sowie dauerhaft
geistig Gestörte. Ihre Willenserklärungen sind nichtig (§\,105 BGB) --- sie können
gar keine wirksamen Verträge schließen.

\textbf{Beschränkt geschäftsfähig} sind Minderjährige zwischen 7 und 17~Jahren
(§\,106 BGB). Sie können wirksam nur lediglich rechtlich vorteilhafte Geschäfte
selbständig abschließen (§\,107 BGB) --- also z.\,B.\ ein Geschenk annehmen. Für
Verträge mit Verpflichtung (etwa Kauf, Miete) brauchen sie die Zustimmung der Eltern.

\textbf{Voll geschäftsfähig} sind Personen ab dem 18.~Geburtstag --- sie können alle
Rechtsgeschäfte ohne Einschränkung abschließen.

\begin{merksatz}[title={§\,110 BGB Taschengeldparagraph}]
Ein vom Minderjährigen ohne Zustimmung des gesetzlichen Vertreters geschlossener Vertrag
gilt als \emph{von Anfang an wirksam}, wenn der Minderjährige die vertragsmäßige
Leistung mit Mitteln bewirkt, die ihm zur freien Verfügung überlassen wurden
(z.\,B.\ Taschengeld). Das ist der berühmte ,,Taschengeld\-paragraph''.
\end{merksatz}

\begin{beispiel}
Ein 6-jähriges Kind kauft im Laden einen Kaugummi für 50~Cent. \emph{Ergebnis:} Der
Kauf ist nichtig, weil das Kind geschäftsunfähig ist. Das Kind hat einen Rückgabeanspruch.

Ein 16-jähriger Auszubildender bestellt online ein Tablet für 400~€. \emph{Ergebnis:}
Der Vertrag ist zunächst schwebend unwirksam --- die Eltern müssen zustimmen. Geschieht
das nicht, ist der Vertrag endgültig unwirksam (§§\,107, 108 BGB). Anders wäre es nur,
wenn der Auszubildende die 400~€ aus seinem freiverfügbaren Taschengeld zahlt
(§\,110 BGB).
\end{beispiel}

\section{Rechtsobjekte}

Während Personen die Träger der Rechte sind, sind \emph{Rechtsobjekte} das, worüber
diese Rechte sich erstrecken. Man unterscheidet:

\begin{itemize}[itemsep=0pt]
  \item \textbf{Sachen}: bewegliche (Auto, Aktenordner, Computer) und unbewegliche (Grundstücke, Gebäude). Innerhalb der beweglichen unterscheidet man weiter zwischen vertretbaren Sachen (austauschbar nach Maß, Zahl oder Gewicht --- z.\,B.\ Druckerpapier) und unvertretbaren Sachen (einmalig --- z.\,B.\ ein Original-Gemälde).
  \item \textbf{Rechte und Forderungen}: Patente, Markenrechte, Urheberrechte, Forderungen aus Verträgen.
\end{itemize}

\section{Arten von Rechtsgeschäften}

\subsection{Einseitig oder mehrseitig}

\begin{merksatz}
Ein \textbf{Rechtsgeschäft} ist eine private Willenserklärung, die einen rechtlichen
Erfolg herbeiführen soll. Es entsteht durch eine (einseitiges) oder mehrere übereinstimmende
(mehrseitiges) Willenserklärung(en).
\end{merksatz}

Ein \emph{einseitiges Rechtsgeschäft} braucht nur die Willenserklärung einer Person.
Beispiele: Kündigung, Testament, Anfechtung, Auslobung. Manche einseitige Rechtsgeschäfte
sind \emph{empfangs­bedürftig} (die Erklärung muss dem anderen zugehen, z.\,B.\ Kündigung),
andere nicht (z.\,B.\ Testament).

Ein \emph{mehrseitiges Rechtsgeschäft} ist ein \emph{Vertrag} --- es braucht mindestens
zwei übereinstimmende Willenserklärungen. Beispiele: Kaufvertrag, Mietvertrag,
Arbeitsvertrag.

\subsection{Verpflichtungs- und Verfügungsgeschäft}

Hier wird es spannend: Im deutschen Recht gilt das \emph{Trennungs- und Abstraktions­
prinzip}. Beim Kauf eines Aktenordners passieren juristisch zwei rechtlich getrennte
Vorgänge.

Das \emph{Verpflichtungs­geschäft} ist der Kaufvertrag selbst (§\,433 BGB). Er begründet
die gegenseitigen Verpflichtungen: Der Verkäufer verpflichtet sich zur Übergabe und
Eigentums­übertragung, der Käufer zur Zahlung des Kaufpreises.

Das \emph{Verfügungs­geschäft} ist die tatsächliche Erfüllung dieser Verpflichtung:
Übergabe der Sache und Eigentumsübertragung (§\,929 BGB) sowie Geldzahlung.

\begin{eselsbruecke}
\textbf{,,Versprechen --- Erfüllen.''} Erst verspreche ich (Verpflichtungs­geschäft =
Vertrag), dann erfülle ich (Verfügungs­geschäft = Übergabe). Beide sind rechtlich
selbständig --- ein Mangel des einen macht das andere nicht automatisch unwirksam
(\emph{Abstraktions­prinzip}).
\end{eselsbruecke}

\section{Vertragsfreiheit und ihre Grenzen}

\subsection{Was Sie alles frei vereinbaren dürfen}

Im privaten Recht gilt der Grundsatz der \emph{Vertragsfreiheit}. Sie umfasst:

\begin{itemize}[itemsep=0pt]
  \item \textbf{Abschlussfreiheit} --- ob und mit wem Sie Verträge schließen
  \item \textbf{Inhaltsfreiheit} --- was im Vertrag steht
  \item \textbf{Formfreiheit} --- mündlich, schriftlich, per Handschlag (Ausnahmen: Grundstückskaufvertrag muss notariell, §\,311b BGB)
  \item \textbf{Aufhebungsfreiheit} --- einvernehmliche Auflösung
\end{itemize}

\subsection{Wo der Gesetzgeber die Freiheit begrenzt}

Die Freiheit endet dort, wo wichtige Rechtsgüter geschützt werden müssen.

\textbf{§\,134 BGB --- Verbotsgesetz}: Verträge, die gegen ein gesetzliches Verbot
verstoßen, sind nichtig. Beispiel: Kaufvertrag über Drogen, Schwarzarbeitsvertrag.

\textbf{§\,138 BGB --- Sittenwidrigkeit und Wucher}: Verträge, die gegen die guten
Sitten verstoßen, sind nichtig. Wucher (Ausnutzen einer Notlage zu auffälligem
Missverhältnis von Leistung und Gegenleistung) ist ein Spezialfall der Sittenwidrigkeit.
Beispiel: Kreditvertrag mit 200\,\% Zinsen p.\,a.\ wegen Notlage des Schuldners.

Daneben gibt es \emph{Form­vorschriften} (Grundstücks­kauf, Eheschließung) und ---
ausnahmsweise --- einen \emph{Kontrahierungs­zwang} (z.\,B.\ Strom­versorger müssen
Bürger versorgen, Krankenhäuser müssen Notfälle aufnehmen).

\section{Zustandekommen eines Kaufvertrags}

\subsection{Antrag und Annahme --- die Grundregel}

\begin{merksatz}
Ein Kaufvertrag kommt zustande durch zwei \emph{übereinstimmende, sich aufeinander
beziehende} Willenserklärungen: \textbf{Antrag (Angebot)} und \textbf{Annahme}.
\end{merksatz}

Der Antrag (\S\,145 BGB) ist die Willenserklärung, mit der jemand einem anderen die
Schließung eines Vertrags anträgt. Der Antragende ist daran gebunden --- er kann den
Antrag nicht beliebig widerrufen, solange die Annahmefrist läuft (§§\,147--149 BGB).

Damit der Antrag wirksam ist, muss er die wesentlichen Vertragspunkte enthalten ---
die sogenannten \emph{essentialia negotii}:

\begin{enumerate}
  \item \textbf{Vertragspartner} (wer ist Käufer, wer Verkäufer)
  \item \textbf{Ware bzw.\ Leistung} (was wird gekauft)
  \item \textbf{Preis} (wie viel)
\end{enumerate}

Die \emph{Annahme} muss vorbehaltlos und rechtzeitig erfolgen. Eine Zustimmung mit
Änderungen gilt als Ablehnung verbunden mit einem neuen Antrag (§\,150 Abs.\,2 BGB).

\begin{beispiel}
Lieferant A bietet 100~Aktenordner für 422\,€ an. Sie schreiben zurück: ,,Wir nehmen
das Angebot an, aber bitte für 400\,€.'' --- Das ist \emph{keine Annahme}, sondern ein
neuer Antrag. Lieferant A muss nun zustimmen, sonst kommt kein Vertrag zustande.
\end{beispiel}

\subsection{Vertragsschluss in der Praxis: Onlineshop}

Frau Meier sieht im Onlineshop von \emph{Bürofix AG} Aktenordner für 4{,}90~€ pro Stück.
Wie kommt der Vertrag zustande?

\begin{enumerate}
  \item Die Anzeige im Onlineshop ist eine \emph{invitatio ad offerendum} --- noch
    \textbf{kein} bindendes Angebot. Der Shop kann Bestellungen ablehnen
    (z.\,B.\ wenn die Ware nicht mehr verfügbar ist).
  \item Frau Meier sendet die Bestellung ab. Damit macht sie den \textbf{Antrag}.
  \item Der Shop sendet eine \textbf{Bestellbestätigung} --- erst diese ist die
    Annahme. Jetzt ist der Vertrag geschlossen.
\end{enumerate}

Achtung: Eine automatische ,,Eingangs­bestätigung'' direkt nach Absenden der Bestellung
ist meist \emph{nicht} die Annahme, sondern nur eine Bestätigung, dass die Bestellung
beim Shop angekommen ist. Die rechtlich verbindliche Annahme erfolgt mit der
,,Versand­bestätigung'' oder einer ausdrücklichen Auftrags­bestätigung.

\subsection{Pflichten aus dem Kaufvertrag (§\,433 BGB)}

\begin{merksatz}[title={§\,433 BGB --- die Hauptpflichten}]
\textbf{Verkäufer:} Sache übergeben \emph{und} Eigentum verschaffen.\\
\textbf{Käufer:} Kaufpreis zahlen \emph{und} Sache abnehmen.
\end{merksatz}

Daneben bestehen \emph{Nebenpflichten}: Information, Aufklärung, Schutz, Mitwirkung.
Die Verletzung von Nebenpflichten kann zu Schadensersatzansprüchen führen.

\section{Allgemeine Geschäftsbedingungen (AGB)}

\subsection{Was sind AGB?}

\begin{merksatz}
\textbf{AGB} sind für eine Vielzahl von Verträgen vorformulierte Vertrags­bedingungen,
die eine Vertragspartei (der \emph{Verwender}) der anderen Partei bei Vertragsschluss
stellt (§\,305 Abs.\,1 BGB).
\end{merksatz}

Beispiele: Onlineshop-AGB, Versicherungs-AGB, Mietvertrags-AGB, Lieferanten-AGB im B2B.
AGB sparen Verhandlungs­zeit, weil nicht jeder Vertrag neu ausgehandelt werden muss.
Die Kehrseite: Die andere Partei kann meist nur ,,akzeptieren oder ablehnen''.
Deshalb schützt der Gesetzgeber sie --- besonders Verbraucher.

\subsection{Wirksame Einbeziehung (§\,305 Abs.\,2 BGB)}

AGB werden nur dann Vertragsbestandteil, wenn drei Voraussetzungen erfüllt sind:

\begin{enumerate}
  \item Der Verwender weist die andere Partei \emph{vor oder bei} Vertragsschluss auf die AGB hin.
  \item Die andere Partei hat die Möglichkeit zur \emph{zumutbaren Kenntnisnahme} (z.\,B.\ Aushang, Link, Beilage).
  \item Die andere Partei \emph{erklärt sich einverstanden} (ausdrücklich oder konkludent).
\end{enumerate}

Im B2B-Geschäft sind die Anforderungen abgeschwächt: Branchen­übliche AGB können auch
ohne ausdrücklichen Hinweis gelten, wenn beide Parteien Kaufleute sind.

\subsection{Überraschende Klauseln und Inhaltskontrolle}

\textbf{§\,305c BGB}: Klauseln, die nach den Umständen so ungewöhnlich sind, dass die
andere Partei nicht damit rechnen musste, werden nicht Vertragsbestandteil. Beispiel:
Eine Allgefahren­versicherung enthält versteckt eine ,,Kunden­abfindungs­klausel'' für
Kosten der Akquise --- so eine Klausel ist überraschend und damit unwirksam.

\textbf{§\,307 BGB --- Inhaltskontrolle}: Klauseln, die den Vertragspartner unangemessen
benachteiligen, sind unwirksam. An ihre Stelle tritt das gesetzliche Recht.

\begin{beispiel}
In den AGB der \emph{Bürofix AG} steht: ,,Reklamationen müssen innerhalb von 24~Stunden
nach Lieferung schriftlich angezeigt werden, sonst gilt die Lieferung als angenommen.''

Diese Klausel verkürzt die gesetzlichen Mängelrechte (§\,437 BGB --- Verjährung in der
Regel zwei Jahre) drastisch und benachteiligt den Käufer unangemessen. \emph{Folge:}
Die Klausel ist nach §\,307 BGB unwirksam. Stattdessen gilt das Gesetz --- der Käufer
hat die volle gesetzliche Verjährungsfrist für Mängelrügen.
\end{beispiel}

\subsection{Kollidierende AGB}

Was, wenn beide Parteien eigene AGB haben, die sich widersprechen? Beispiel: Bürofix
schreibt ,,Erfüllungsort München'', Sie schreiben in Ihren Einkaufs-AGB ,,Erfüllungs­
ort Berlin''. \emph{Lösung:} Bei kollidierenden AGB werden die widersprüchlichen
Klauseln nicht Vertragsbestandteil. Stattdessen greift das gesetzliche Recht
(BGB/HGB). Bei Erfüllungsort: §\,269 BGB --- in der Regel der Wohnsitz/Sitz des
Schuldners (also des Verkäufers).

\section{Zusammenfassung}

\begin{enumerate}
  \item \textbf{Privates Recht} (BGB, HGB) regelt Beschaffungsverträge.
  \item \textbf{Geschäftsfähigkeit}: 0--6 (unfähig), 7--17 (beschränkt), ab 18 (voll); §\,110 Taschengeld\-paragraph als wichtige Ausnahme.
  \item \textbf{Rechtsgeschäfte}: einseitig (Kündigung) oder mehrseitig (Vertrag). Trennungsprinzip: Verpflichtung $\neq$ Verfügung.
  \item \textbf{Vertragsfreiheit} mit Grenzen (§§\,134, 138 BGB).
  \item \textbf{Kaufvertrag}: durch Antrag (§\,145) + Annahme; Pflichten in §\,433 BGB.
  \item \textbf{AGB}: Einbeziehung nach §\,305 Abs.\,2; Inhaltskontrolle nach §\,307 BGB.
\end{enumerate}

\newpage

\section*{Glossar}
\addcontentsline{toc}{section}{Glossar}

\glossareintrag{Privates Recht}{Rechtsgebiet, das Beziehungen zwischen gleichberechtigten Personen regelt (BGB, HGB).}

\glossareintrag{Öffentliches Recht}{Rechtsgebiet, das das Verhältnis zwischen Bürger und Staat regelt (Grundgesetz, Verwaltungsrecht).}

\glossareintrag{Natürliche Person}{Jeder Mensch ab Geburt; rechtsfähig nach §\,1 BGB.}

\glossareintrag{Juristische Person}{Künstliche Rechtspersönlichkeit (GmbH, AG, e.\,V.\, Stadt) --- selbst Träger von Rechten und Pflichten.}

\glossareintrag{Geschäftsfähigkeit}{Fähigkeit, eigenständig wirksame Rechtsgeschäfte vorzunehmen. Drei Stufen: unfähig, beschränkt, voll.}

\glossareintrag{Beschränkt geschäftsfähig}{Person zwischen 7 und 17~Jahren --- nur lediglich rechtlich vorteilhafte Geschäfte selbständig wirksam (§\,107 BGB).}

\glossareintrag{Taschengeldparagraph}{§\,110 BGB --- Vertrag eines Minderjährigen ist wirksam, wenn die Leistung mit frei verfügbaren Mitteln bewirkt wird.}

\glossareintrag{Rechtsgeschäft}{Private Willenserklärung mit dem Ziel, einen rechtlichen Erfolg herbeizuführen.}

\glossareintrag{Einseitiges Rechtsgeschäft}{Nur eine Willenserklärung nötig; z.\,B.\ Kündigung, Testament.}

\glossareintrag{Mehrseitiges Rechtsgeschäft (Vertrag)}{Mindestens zwei übereinstimmende Willenserklärungen; z.\,B.\ Kaufvertrag.}

\glossareintrag{Verpflichtungsgeschäft}{Rechtsgeschäft, das eine Verpflichtung begründet (z.\,B.\ Kaufvertrag).}

\glossareintrag{Verfügungsgeschäft}{Rechtsgeschäft, das eine Verpflichtung erfüllt (z.\,B.\ Übergabe \& Eigentumsübertragung nach §\,929 BGB).}

\glossareintrag{Trennungs- und Abstraktionsprinzip}{Verpflichtung und Verfügung sind rechtlich getrennt und unabhängig voneinander wirksam.}

\glossareintrag{Vertragsfreiheit}{Grundsatz: Jeder darf grundsätzlich frei entscheiden, ob, mit wem und worüber er Verträge schließt.}

\glossareintrag{Verbotsgesetz (§\,134 BGB)}{Vertrag, der gegen ein gesetzliches Verbot verstößt, ist nichtig.}

\glossareintrag{Sittenwidrigkeit (§\,138 BGB)}{Vertrag, der gegen die guten Sitten verstößt, ist nichtig. Wucher = Spezialfall.}

\glossareintrag{Antrag (Angebot)}{Bindender Vertragsantrag nach §\,145 BGB; muss die essentialia negotii enthalten.}

\glossareintrag{Annahme}{Vorbehaltlose Zustimmung zum Antrag; mit Änderungen gilt sie als neuer Antrag (§\,150 Abs.\,2 BGB).}

\glossareintrag{Essentialia negotii}{Wesentliche Vertragsbestandteile --- beim Kaufvertrag: Vertragspartner, Ware, Preis.}

\glossareintrag{Invitatio ad offerendum}{Aufforderung zur Abgabe eines Angebots --- selbst noch kein bindender Antrag (z.\,B.\ Schaufenster, Onlineshop-Anzeige).}

\glossareintrag{§\,433 BGB}{Hauptpflichten beim Kaufvertrag: Verkäufer übergibt + verschafft Eigentum, Käufer zahlt + nimmt ab.}

\glossareintrag{AGB}{Allgemeine Geschäftsbedingungen --- für eine Vielzahl von Verträgen vorformulierte Bedingungen (§\,305 Abs.\,1 BGB).}

\glossareintrag{Überraschende Klausel (§\,305c BGB)}{Ungewöhnliche AGB-Klausel, mit der die andere Partei nicht rechnen musste --- nicht Vertragsbestandteil.}

\glossareintrag{Inhaltskontrolle (§\,307 BGB)}{Prüfung, ob AGB-Klauseln die andere Partei unangemessen benachteiligen --- bei Verstoß: unwirksam.}

\end{document}
